\section{Introduction}

Sign Language Translation (SLT) refers to the process of converting a continuous sign language video into a corresponding written language translation, with the aim of bridging communication gaps between the deaf and hearing communities~\cite{de2023machine}. SLT is a complex task that can be roughly decomposed into three sub-problems: 1) Sign Language Segmentation (SLS): segmenting a video of continuous sign language into many individual signs; 2) Sign Language Recognition (SLR): classifying a video of a single sign; and 3) translation of the identified signs into a written language whose grammar differs from the grammar of the sign language.

This grammatical difference adds a layer of complexity for interpreting SLT models as there is no one-to-one mapping between signs and words, and the ordering might vary. Therefore, it is significantly complex to interpret a SLT model without expert knowledge on both the grammar of the signs and the target written language. Furthermore, there is no universal sign language since each country or region has its own sign language with a particular grammar. Translation from video to Sign Language Glosses -written representation of the signs-~\cite{camgoz2020sign} constitutes an intermediate approach that can help us bypass these difficulties to better understand SLT models.

In latest years, SLT models have been improved by new deep learning architectures like the Transformer networks~\cite{camgoz2020sign,yin2020better,de2021frozen,aloysius2021incorporating}, which rely mainly on the attention mechanism. This mechanism has been proven to be interpretable and directly correlated with feature importance for tasks such as Natural Language Translation \cite{vashishth2019attention}.

% The complexity of the task, alongside its multi modal and temporal nature invites to explore both the feasibility of existing methods as the development of new ones that consider the peculiarities of the field. 

In this work, we present, to the best of our knowledge, the first comprehensive analysis of interpretability of a Sign Language Translation Model, by training and studying a Transformer architecture for the Greek Sign Language using the well-known Greek Sign Language Dataset~\cite{adaloglou2020comprehensive}. We  study the attention outputs within the decoder at each decoding step, which consists of the additive combination of \textit{cross-attention}, performed over the context vector generated by the encoder, and \textit{self-attention}, that operates between tokens of the target sentence. By examining these outputs, we gain insight into how the model allocates attention across different frames and target words, thus identifying which elements are most significant for SLT based on their attention scores and magnitudes. Our work not only provides new insights into the relationship between written and sign languages but also paves the way for developing more transparent and explainable translation systems, which are crucial for deploying SLT technology in real-world applications where understanding model decisions is essential for both accuracy and user trust.